% Flavor catalog per-process template.
% process_id: CR012
% family: collider_rs
% owner_agent_id: PKA-CR012
% writer_agent_id: null
% checker_agent_id: null
% opus_signoff_id: null
% Canonical metadata sidecar: CR012.yaml
% Status is controlled by the YAML sidecar status_history, not this comment.

\section{\texorpdfstring{Diboson High-Mass Resonance, \(V'\to VV\)}{Diboson High-Mass Resonance, Vprime to VV}}

\subsection*{Process}
This entry covers direct LHC searches for a high-mass spin-1 resonance in
diboson final states,
\[
  pp\to (W',Z',G'_{\rm vector}) \to WW,\;WZ,\;ZZ,
\]
with the phenomenological interpretation usually expressed in the heavy-vector
triplet (HVT) language.  The \(G'_{\rm vector}\) label is treated here as a
generic heavy gauge-resonance placeholder; the CMS bulk-graviton rows are useful
adjacent RS context but are spin-2 and are not the canonical CR012 spin-1 value.

\subsection*{Current best limits}
The canonical current spin-1 diboson number is the CMS full-Run-2 all-jets
analysis \texttt{CMS-B2G-20-009} / arXiv:2210.00043, cross-checked against the
PDG 2025 heavy-boson listing.  For HVT model B with \(g_V=3\), CMS excludes an
HVT \(W'\to WZ\) resonance below \(4.4~\mathrm{TeV}\) at \(95\%\) CL
(\texttt{pdg2025\_heavy\_bosons\_wprime\_wz.txt};
\texttt{cms\_2023\_b2g\_20\_009\_arxiv2210\_00043.txt}).  The same CMS Table 2
gives a mass-degenerate \(V'\to VV\) lower limit of \(4.5~\mathrm{TeV}\), a
stronger \(4.8~\mathrm{TeV}\) limit only after combining \(VV+VH\), and a
disjoint observed \(Z'\to WW\) exclusion in the intervals
\(1.3\)--\(3.1~\mathrm{TeV}\) and \(3.3\)--\(3.5~\mathrm{TeV}\).  The
\(4.8~\mathrm{TeV}\) number should therefore be quoted as a degenerate-HVT
\(VV+VH\) result, not as a pure-diboson-only bound.

\subsection*{Relevance to RS / anarchic-flavor pipeline}
Custodial and composite RS models generically contain electroweak KK vectors
whose longitudinal \(W/Z\) branching fractions can be large.  The HVT model-B
benchmark is not an RS spectrum calculation, but it is a standard simplified
proxy for strongly coupled vector triplets and is the cleanest public
collider-language bridge for \(W'\) and \(Z'\) diboson searches.

For the anarchic-flavor pipeline, this is not the leading scale setter.  The
internal quark-scan methodology note quotes
\(M_{KK}^{\min}(p50,g_{s*}=3)=47.26~\mathrm{TeV}\) for the low-energy
\(\Delta F=2\)-dominated scan
(\texttt{internal\_quark\_scan\_methodology\_extract.txt}).  A direct LHC
diboson reach of order \(4\)--\(5~\mathrm{TeV}\) is therefore a cross-check and
a handle on non-anarchic, custodial, or otherwise flavor-aligned RS variants,
rather than a competing bound for the default anarchic scan.

\subsection*{Post-2010 developments}
Pappadopulo--Thamm--Torre--Wulzer, arXiv:1402.4431, introduced the HVT
framework used by ATLAS and CMS to translate \(VV\) and \(VH\) searches into a
small set of vector-resonance parameters.  ATLAS arXiv:1506.00962 then made the
boosted all-hadronic \(WW/WZ/ZZ\) search a central LHC channel and recorded the
well-known 8 TeV \(WZ\)-region fluctuation.  CMS arXiv:1705.09171 provided the
first combined \(VV/VH\) HVT interpretation across 8 and early 13 TeV searches.
The 13 TeV all-hadronic searches in ATLAS arXiv:1708.04445 and CMS
arXiv:1708.05379 moved the channel into the multi-TeV regime with boosted
boson-tagging.

The 2019 full-Run-2/partial-Run-2 updates, ATLAS arXiv:1906.08589 and CMS
arXiv:1906.05977, established the modern high-mass all-jets analyses.  CMS
arXiv:2109.06055 added an independent semileptonic Run-2 channel for \(WW\),
\(WZ\), and \(WH\).  The current all-jets CMS arXiv:2210.00043 result then used
the full Run-2 data set, machine-learning boson taggers, and simultaneous
\(VV/VH\) categories to give the strongest public HVT spin-1 statement used
above.

\subsection*{Constraint validity and model dependence}
These are direct-search limits, not model-independent limits on the RS
Kaluza--Klein scale.  The HVT interpretation assumes a narrow resonance, a
specific relation between fermion and boson couplings, and production through
the light-quark or vector-boson initial states present in the experimental
templates.  In a concrete RS model the mapping depends on the custodial gauge
group, fermion embeddings, brane kinetic terms, total width, and branching
fractions to \(t\bar t\), \(VH\), leptons, dijets, or invisible/exotic modes.

The distinction between \(V'\to VV\) and \(V'\to VV+VH\) is load-bearing for
CR012.  A model with a large \(VH\) branching fraction can use the combined CMS
\(4.8~\mathrm{TeV}\) number; a diboson-only reinterpretation should instead use
the \(W'\to WZ\), \(Z'\to WW\), or \(V'\to VV\) rows and their corresponding
assumptions.

\subsection*{Code coverage in this repo}
\textbf{NO}.  Greps over \texttt{quarkConstraints/}, \texttt{qcd/},
\texttt{flavorConstraints/}, \texttt{neutrinos/}, \texttt{yukawa/},
\texttt{warpConfig/}, \texttt{solvers/}, \texttt{scanParams/}, and
\texttt{tests/} found no diboson, HVT, ATLAS/CMS search, HEPData likelihood, or
reinterpretation-tool implementation.  The only adjacent matches are KK-scale
bookkeeping and low-energy flavor machinery: \texttt{quarkConstraints/scan.py:359}
maps \(\Lambda_{\rm IR}\) to \(M_{KK}\), \texttt{quarkConstraints/scan.py:377}
calls \(\Delta F=2\) constraints, \texttt{quarkConstraints/deltaf2.py:320}
documents tree-level KK-gluon-inspired \(\Delta F=2\) Wilsons, and
\texttt{scanParams/scan.py:523} runs the unrelated \(\mu\to e\gamma\) LFV
filter.  The only ``CMS'' grep hit is \texttt{tests/test\_alpha\_s.py:89}, a
RunDec QCD-running example, not a collider search.

\subsection*{Implementation difficulty}
\textbf{HIGH}.  A faithful implementation needs either collaboration likelihood
material or an external reinterpretation chain such as CheckMATE, MadAnalysis5,
or SModelS, plus an RS spectrum/branching-ratio calculation.  A single scalar
mass threshold would be misleading because the limit changes with width,
production mode, \(VV/VH\) branching fractions, and custodial embeddings.

\subsection*{Key references}
\texttt{PDG2025\_HeavyBosons\_Wprime\_WZ};
\texttt{CMS2023\_B2G\_20\_009};
\texttt{HEPData2022\_CMS\_B2G\_20\_009};
\texttt{ATLAS2019\_Diboson139fb};
\texttt{CMS2019\_B2G\_18\_002};
\texttt{CMS2018\_Dijet\_B2G\_17\_001};
\texttt{ATLAS2017\_Diboson36fb};
\texttt{CMS2017\_BosonCombination};
\texttt{ATLAS2015\_8TeVDiboson};
\texttt{CMS2022\_Semileptonic\_B2G\_19\_002};
\texttt{Pappadopulo2014\_HVT};
\texttt{InternalQuarkScanMethodology\_RUNA\_47TeV}.
